\begin{abstract}
Standard convolutional neural networks (CNNs) exploit translational symmetry but remain
blind to rotational structure, typically compensating through data augmentation at training
time. Group-equivariant convolutional neural networks (\gcnn{}s) instead encode rotational and
reflective symmetry directly into the network architecture, guaranteeing equivariance by
construction at every layer. Histopathology images are a natural domain for this inductive bias:
tissue slides are acquired without a canonical orientation, making rotational symmetry a strong 
and natural inductive bias, rather than an approximation.

This work studies whether the structural symmetry of \gcnn{}s translates into measurable
benefits for colorectal tissue type classification on the NCT-CRC-HE-100K dataset, a
nine-class benchmark covering the major tissue compartments found in human colorectal cancer
slides. Three conditions are compared: (i) a standard CNN baseline, (ii) the same CNN trained
with rotation and reflection augmentation, and (iii) a \gcnn{} using the \(p4m\) symmetry group.
Critically, we vary the fraction of available training data to probe \emph{sample efficiency}:
the hypothesis is that built-in equivariance should confer the largest advantage when labelled
data is scarce, where augmentation alone cannot substitute for a principled structural prior.
Our experiments are implemented in PyTorch and managed with Hydra for reproducible
configuration. We report classification accuracy, macro-averaged F1 score, and per-class
results across training-set fractions ranging from 5\% to 100\%.
\end{abstract}